% Consequence of the single Brauer-zero projective-count reduction.
\begin{corollary}[Proved: simple-product quotients selected by Brauer-zero representations]
\label{D:Brauer-zero-selected-simple-product}
Let $Q$ be a group in one of the global or complete-local cases
of Theorem~\ref{D:Brauer-zero-projective-representation-growth},
Corollaries~\ref{D:Brauer-zero-detecting-surface-transfer} and
\ref{D:Brauer-zero-normalization-transfer}, or
Theorem~\ref{D:local-Brauer-zero-projective-growth}.
Use exactly the Brauer-vanishing condition of the chosen statement.
Fix $\epsilon>0$. For each isomorphism type $H$ in a collection
of nonabelian finite simple groups, choose one faithful embedding
\[
       \iota_H:H\hookrightarrow
                    \operatorname{PGL}_{r_H}(\mathbf F_{q_H}),
       \qquad q_H=p^{f_H},\qquad
                    |H|\ge q_H^{\epsilon r_H^2}.
\]
Let $\mathcal K_H$ consist of the distinct kernels of continuous
epimorphisms $\eta:Q\twoheadrightarrow H$ for which
$\iota_H\eta$ satisfies that Brauer-vanishing condition.
Then
\[
           Q\big/\bigcap_H\bigcap_{K\in\mathcal K_H}K
\]
is topologically finitely generated. Its simple-factor
multiplicities satisfy $|\mathcal K_H|\le |H|^{C/\epsilon}$
for one constant $C$ independent of $H$.
Only the selected kernels are included; no assertion about
the other kernels of the same simple isomorphism types is made.
\end{corollary}
\begin{proof}
The relevant Brauer-zero theorem bounds the ambient projective
conjugacy classes by $q_H^{Cr_H^2}$. Distinct epimorphism kernels
give distinct such classes, because conjugation does not change
a kernel. The same uniform $C$ works for all the chosen fields
and ranks, so
\[
             |\mathcal K_H|\le q_H^{Cr_H^2}
                                  \le |H|^{C/\epsilon}.
\]
For precision, put $a=C/\epsilon$ and
$d=\lceil a\rceil+3$. Equations(2.1)--(2.2) in the proof of
Proposition2.6(ii), printed pp12--13, of
\citet{LucchiniMarionTracey2020D}, with its constant
$\gamma\ge1/2$, give
\[
 \psi_H(d)\ge\frac{|H|^{d-1}}{2|\operatorname{Out}(H)|}
              \ge\frac{|H|^{d-2}}2\ge |H|^a.
\]
Here $\psi_H(d)$ is the largest multiplicity of a direct power
of $H$ generated by $d$ elements. We used the outer-automorphism
bound on p13 of the same source, in particular
$|\operatorname{Out}(H)|\le |H|$, and $d(H)=2$.
Thus every $H^{|\mathcal K_H|}$ is generated by $d$ elements.
Products of different simple isomorphism types have no common
nontrivial simple quotients, so combining these tuples still
generates every finite subproduct, by Goursat's lemma.

Goursat's lemma identifies the image of $Q$ on every finite
subcollection of distinct simple quotient kernels with the full
product. Indeed a proper subdirect product of nonabelian simple
factors identifies two isomorphic simple quotients, and hence
identifies their kernels, contrary to their selection as distinct.
Thus the displayed quotient is the full profinite product of
the selected simple factors. The uniform finite-product generator
bound gives a generating tuple for this product by compactness:
in a fixed power of the product, the closed conditions that a
tuple generate each finite subproduct have the finite intersection
property. This proves the conclusion.
\end{proof}

\begin{remark}[Exact limitation]
This is a direct consequence of the Brauer-zero count, not an
additional movement event. The remaining projective simple
quotient kernels can have nonzero geometric Brauer class on the
relevant surface or local puncture. General monolithic crowns,
alternating multiplicities and the other uncovered families
remain separate.
\end{remark}
